\documentclass[11pt]{article}
\usepackage{amsmath,amssymb,amsthm}
\usepackage{geometry}
\geometry{margin=1in}
\usepackage{hyperref}
\usepackage{booktabs}
\usepackage{enumerate}
\usepackage{microtype}
\usepackage{xcolor}

% Theorem environments
\newtheorem{theorem}{Theorem}[section]
\newtheorem{proposition}[theorem]{Proposition}
\newtheorem{corollary}[theorem]{Corollary}
\newtheorem{lemma}[theorem]{Lemma}
\theoremstyle{definition}
\newtheorem{definition}[theorem]{Definition}
\theoremstyle{remark}
\newtheorem{remark}[theorem]{Remark}

% Convenience macros
\newcommand{\hqa}{H_{\mathrm{QA}}}
\newcommand{\hraw}{H_{\mathrm{raw}}}
\newcommand{\etaeff}{\eta_{\mathrm{eff}}}
\newcommand{\kmin}{\kappa_{\min}}
\newcommand{\rhoO}{\rho(\mathcal{O})}
\newcommand{\rhoPL}{\rho_{\mathrm{PL}}(\mathcal{O})}
\newcommand{\Zfi}{\mathbb{Z}[\varphi]}
\newcommand{\Qfi}{\mathbb{Q}(\sqrt{5})}
% \pmod provided by amsmath

\title{\textbf{Unified Curvature Certificates Across Learning, Aggregation,\\
Attention, Arithmetic, and Symbolic Search}}
\author{[Author redacted for review]}
\date{}

\begin{document}
\maketitle

\begin{abstract}
The dynamics of the Quantum Arithmetic (QA) system correspond to multiplication
by $\varphi^2$ in $\Zfi$, the ring of integers of $\Qfi$: the map
$T:(b,e)\mapsto(b{+}e,\,b{+}2e)$ equals $Q^2$, the square of the Fibonacci
companion matrix, and acts as ${\times}\varphi^2$ on $b+e\varphi$. The Fibonacci
quadratic norm $N(b+e\varphi)=b^2+be-e^2$ is $T$-invariant, orbit lengths equal
$\pi(m)/2$ (half the Pisano period), and the harmonic index $\hqa$ admits a clean
operator decomposition: $4\hraw=\mathrm{tr}(\Lambda)$ where
$\Lambda=\mathrm{diag}(\sigma_d,\sigma_{od})$ is a $2\times 2$ cross-coupling
operator whose degree-0 eigenvalue $\sigma_d=ba/(e^2+d^2)$ tracks the projective
Fibonacci angle and whose degree-1 eigenvalue $\sigma_{od}=ed/(b+a)$ collapses to
$e/2$ via the Fibonacci identity $a+b=2d$.

Building on this algebraic foundation, we introduce a certificate-normal-form
framework for comparing one-step stability across heterogeneous computational systems.
For any certified family whose local update writes as
$p_{\mathrm{after}}=p_{\mathrm{before}}-\etaeff\cdot\mathrm{grad}$ with
$\etaeff=\mathrm{lr}\cdot\mathrm{gain}\cdot\hqa$, we define the universal
curvature score
\[
  \kappa = 1 - |1-\etaeff|.
\]
We instantiate this across nine certified families: gradient-based learning [89],
graph aggregation [93], attention layers [94], modular arithmetic dynamics [95],
symbolic search [96], orbit curvature [97], spectral-gain extensions [98,\,99],
and gradient Lipschitz gain [101]. Three families ([98], [99], [101]) derive gain
from native structural objects---the spectral norm of weight/score matrices and the
gradient $\ell_2$ norm---upgrading from consistency check to structural analysis.
For QA orbit families, the finite enumerable orbit structure enables exact multi-step
descent results: if $\kmin(\mathcal{O})>0$, then
$L_{t+L}\leq\rhoO\cdot L_t$ where $\rhoO=\prod_t(1-\kappa_t)^2<1$ is computable
exactly (scalar quadratic loss, Theorem~\ref{thm:descent-quad}); for any
$\beta$-smooth $\mu$-PL loss,
$L(w_{t+L})-L^*\leq\rhoPL\cdot(L(w_t)-L^*)<L(w_t)-L^*$
(Theorem~\ref{thm:descent-pl}), with $\rhoPL=\rhoO$ when $\beta=\mu$.
Empirically: across seven QA substrates, mean $\kappa$ correlates with final loss at
$r=-0.843$; normalizing $\etaeff=1$ across all substrates equalizes convergence
(loss std $1.9\times10^{-5}$), confirming that $\hqa$ governs convergence
exclusively through $\etaeff$.
\end{abstract}

\section{Introduction}

The field of artificial intelligence is characterized by a proliferation of
architectures. From multi-layer perceptrons to Graph Neural Networks (GNNs) and
Transformers, each model family has its own stability considerations and hyperparameter
conventions. Comparing the stability of a GNN's aggregation step with that of an
attention head in a Transformer is not straightforward---each family uses different
update semantics, different scale conventions, and different metadata. However, once
each is mapped to a certified scalar local update form, comparison becomes possible.

This paper addresses this challenge by proposing a unified framework for certifying
one-step stability across diverse computational architectures. Our primary
contribution is the introduction of a curvature score $\kappa$, derived from a
common mathematical substrate $\hqa$. This metric provides a single, interpretable
value characterizing the single-step convergence behavior of a system's certified
update rule, independent of domain-specific implementation details.

\paragraph{What this paper claims.}
The paper's claim is that nine heterogeneous architecture families admit the same
machine-checkable one-step curvature certificate form, not that they share the same
full dynamics. For three of these families ([98], [99], [101]) the gain is not a
free witness but a derived structural invariant---the spectral norm of a native
operator---making the certificate a structural analysis, not merely a consistency
check. For QA orbit families specifically, Theorems~\ref{thm:descent-quad}
and~\ref{thm:descent-pl} provide exact multi-step descent bounds for scalar
quadratic and $\beta$-smooth $\mu$-PL losses respectively. We do not assert general
global convergence, multi-step equivalence across architectures, or semantic identity
across architecture classes.

We demonstrate this across nine certified families:
\begin{enumerate}
  \item \textbf{QALM Gradient [89]:} Standard gradient-based optimization.
  \item \textbf{GNN Aggregation [93]:} The neighborhood aggregation step in GNNs.
  \item \textbf{Attention Layer [94]:} The core mechanism in Transformer models.
  \item \textbf{QARM Arithmetic [95]:} Operations within a modular arithmetic system.
  \item \textbf{Symbolic Search [96]:} Heuristic-guided search algorithms.
  \item \textbf{Orbit Curvature [97]:} QA orbit stability margin across a full finite orbit.
  \item \textbf{GNN Spectral Gain [98]:} Gain derived from $\sigma_{\max}(W)$.
  \item \textbf{Attention Spectral Gain [99]:} Gain derived from $\sigma_{\max}(QK^\top/\!\sqrt{d_k})$.
  \item \textbf{Gradient Lipschitz Gain [101]:} Gain derived from $\|\mathbf{g}\|_2$, capped at 2.0.
\end{enumerate}

\begin{table}[h]
\centering
\caption{Nine certified families at a glance.}
\label{tab:families}
\begin{tabular}{lllll}
\toprule
Family & Domain & Gain witness & Gain type & ID \\
\midrule
QALM Gradient & Gradient optimization & \texttt{gain} & free scalar & [89] \\
GNN Aggregation & Graph neural networks & \texttt{agg\_gain} & free scalar & [93] \\
Attention Layer & Transformer attention & \texttt{attn\_gain} & free scalar & [94] \\
QARM Arithmetic & Modular arithmetic & \texttt{qarm\_gain} & free scalar & [95] \\
Symbolic Search & Heuristic beam search & \texttt{sym\_gain} & free scalar & [96] \\
Orbit Curvature & QA orbit stability & $\kmin$ & derived (orbit enum.) & [97] \\
GNN Spectral Gain & GNN weight geometry & $\sigma_{\max}(W)$ & derived (power iter.) & [98] \\
Attention Spectral Gain & Attention geometry & $\sigma_{\max}(QK^\top/\!\sqrt{d_k})$ & derived (power iter.) & [99] \\
Gradient Lipschitz Gain & Gradient descent & $\min(\|\mathbf{g}\|_2,2)$ & derived (L2 norm) & [101] \\
\bottomrule
\end{tabular}
\end{table}

\section{Mathematical Foundation}
\label{sec:mathematical-foundation}

The $\hqa$ substrate is not an ad-hoc formula. It arises from the arithmetic of
$\Qfi$ and from the projective dynamics of the Fibonacci recurrence. We summarise
the structure here; formal proofs appear in Section~\ref{sec:appendix}.

\paragraph{The ring $\Zfi$ and the QA map.}
The golden ratio $\varphi=(\sqrt{5}+1)/2$ satisfies $\varphi^2=\varphi+1$. The ring
of integers $\Zfi=\{b+e\varphi : b,e\in\mathbb{Z}\}$ carries the Fibonacci quadratic
norm $N(b+e\varphi)=b^2+be-e^2$. The QA map $T:(b,e)\mapsto(d,a)$ with $d=b+e$,
$a=b+2e$ (before modular reduction) corresponds to multiplication by $\varphi^2$:
\[
  \varphi^2(b+e\varphi) = (b+e)+(b+2e)\varphi = d+a\varphi.
\]
The matrix of $T$ is $Q^2=\bigl(\begin{smallmatrix}1&1\\1&2\end{smallmatrix}\bigr)$,
the square of the Fibonacci companion matrix
$Q=\bigl(\begin{smallmatrix}0&1\\1&1\end{smallmatrix}\bigr)$.

\paragraph{Norm invariance and orbit structure.}
Since $N(\varphi^2)=1$, the norm $N(b+e\varphi)=b^2+be-e^2$ is $T$-invariant.
Modulo $m$, the orbit length equals $\pi(m)/2$---half the Pisano period---because
the order of $Q^2$ in $GL_2(\mathbb{Z}/m\mathbb{Z})$ is $\pi(m)/2$ for all $m\geq 3$
(Propositions~\ref{prop:half-pisano} and~\ref{prop:norm-invariant}).
For $m=9$: $\pi(9)=24$, orbit length $=12$, 72 starting pairs ("Cosmos" group).
The modular value $N(b+e\varphi)\bmod m$ classifies orbits: for $m=9$, $3\nmid N$
gives cosmos (length 12), $3^2\mid N$ gives satellite (length 4 or 1), $3^4\mid N$
gives the singularity (length 1).

\paragraph{Projective Fibonacci flow.}
The projective ratio $z=e/(b+e)$ evolves as the M\"obius map $z\mapsto(z+1)/(z+2)$
under $T$, with unique fixed point $z^*=1/\varphi\approx 0.618$. Every pre-modular
orbit converges projectively to $z^*$.

\paragraph{$\hqa$ as a cross-coupling trace.}
The step matrix $M=\bigl(\begin{smallmatrix}b&e\\d&a\end{smallmatrix}\bigr)$
(rows = state before and after $T$) satisfies $\det(M)=N(b+e\varphi)$
(pre-modular). The harmonic index decomposes as:
\[
  4\hraw = \underbrace{\frac{ba}{e^2+d^2}}_{\sigma_d,\;\text{degree-0}}
          + \underbrace{\frac{ed}{b+a}}_{\sigma_{od},\;\text{degree-1}}
          = \mathrm{tr}(\Lambda),
  \qquad \Lambda = \mathrm{diag}(\sigma_d,\,\sigma_{od}).
\]
The degree-0 component $\sigma_d=\cos(2\arctan(z))$ tracks the projective Fibonacci
angle and is scale-invariant; the degree-1 component $\sigma_{od}=e/2$ is fixed by
the Fibonacci identity $a+b=2d$ and carries the amplitude. The factor $1/4$ is the
normalised-trace scaling: $\hraw=\mathrm{tr}(\Lambda)/4=R(\Lambda,(1,1)^\top/\sqrt{2})/2$.
Properties (i) degree-(0,1) homogeneity, (ii) transpose symmetry
$\sigma(M)=\sigma(M^\top)$, and (iii) projective Fibonacci convergence characterise
$4\hraw$ as the natural cross-coupling trace of the QA step matrix
(Section~\ref{sec:appendix}).

\section{The $H_{QA}$ Substrate}
\label{sec:substrate}

Section~\ref{sec:mathematical-foundation} established the algebraic origin of $\hqa$.
We now give the formal computational definition used by all validators.

\subsection{Formal Definition}

Given four positive integers $a, b, d, e \geq 1$, define:
\[
  G = e\cdot e + d\cdot d, \qquad F = b\cdot a.
\]
From these, compute a raw harmonic ratio $\hraw$ with a small $\varepsilon=10^{-12}$:
\[
  \hraw = 0.25 \cdot \left( \frac{F}{G + \varepsilon}
                           + \frac{e \cdot d}{a + b + \varepsilon} \right).
\]

\subsection{Boundedness Property}

The saturating function
\[
  \hqa = \frac{|\hraw|}{1 + |\hraw|}
\]
maps $\hraw$ to $[0,1)$. This ensures $\hqa$ is always non-negative and strictly
less than 1, a critical prerequisite for the stability analysis of the universal
$\kappa$ formula.

\subsection{Geometric Interpretation}

As derived in Section~\ref{sec:mathematical-foundation}, $4\hraw=\mathrm{tr}(\Lambda)$
where $\Lambda=\mathrm{diag}(\sigma_d,\sigma_{od})$ is the cross-coupling operator of
the QA step matrix. The degree-0 eigenvalue $\sigma_d=ba/(e^2+d^2)=\cos(2\theta)$
where $\theta=\arctan(e/d)$; the degree-1 eigenvalue $\sigma_{od}=ed/(b+a)=e/2$
carries the state amplitude. Both terms are algebraically determined by the
$\mathbb{Q}(\sqrt{5})$ structure; the normalisation constant $1/4$ carries a residual
derivation gap (Section~\ref{sec:appendix}).

\section{The Universal $\kappa$ Formula}

\subsection{Derivation}

Consider a general iterative process where a parameter $p$ is updated at each step.
We define an \emph{effective learning rate} modulated by the system's substrate:
\[
  \etaeff = \mathrm{lr} \cdot \mathrm{gain} \cdot \hqa.
\]
Here $\mathrm{lr}$ is a global base learning rate, $\hqa$ is the Harmonic Index of
the system's current state, and $\mathrm{gain}$ is a scalar witness specific to the
architecture class. The certified update rule becomes:
\[
  p_{\mathrm{after}} = p_{\mathrm{before}} - \etaeff \cdot \mathrm{grad}.
\]
We define the curvature score:
\[
  \kappa = 1 - |1-\etaeff| = 1 - |1 - \mathrm{lr}\cdot\mathrm{gain}\cdot\hqa|.
\]
\textbf{Interpretation.} The quantity $\kappa\in(-\infty,1]$ is positive exactly
when $\etaeff\in(0,2)$ and is maximised at $\etaeff=1$. It measures proximity to
the centre of the admissible one-step interval, not the absolute magnitude of the
step.

\subsection{Stability Theorem}

\begin{theorem}[One-Step Stability]
The single-step update is stable ($\kappa>0$) if and only if $\etaeff\in(0,2)$.
\end{theorem}
\begin{proof}
$\kappa>0\Leftrightarrow |1-\etaeff|<1\Leftrightarrow 0<\etaeff<2$. A practical
sufficient condition is $\mathrm{gain}\in(0,2]$ and $\mathrm{lr}<1$; since
$\hqa<1$, this gives $\etaeff<2$.
\end{proof}

\section{Nine Architecture Classes}
\label{sec:families}

We instantiate the framework across nine certified families. Cross-row differences in
Table~\ref{tab:comparison} arise from different substrate and gain inputs, not
architecture identity. Families [98], [99], and [101] derive gain from a native
operator; the remainder supply a free scalar.

\subsection{QALM Gradient [89]}
\textbf{Gain:} \texttt{gain}. \textbf{Description:} The \texttt{gain} parameter
acts as a scalar on the learning rate, modulated by $\hqa$ of the current state.
\textbf{Example:} Substrate $(b,e,d,a)=(1,2,3,5)$, $\hqa\approx 0.2571$.

\subsection{GNN Aggregation [93]}
\textbf{Gain:} \texttt{agg\_gain}. \textbf{Metadata:} \texttt{n\_nodes},
\texttt{n\_edges}. \textbf{Description:} Controls the strength of the aggregated
message updating a node's representation. Structural metadata provides auditable
graph context but does not enter the curvature computation.

\subsection{Attention Layer [94]}
\textbf{Gain:} \texttt{attn\_gain}. \textbf{Metadata:} \texttt{n\_heads},
\texttt{d\_model}, \texttt{seq\_len}. \textbf{Description:} Scales the output of
the attention head.

\subsection{QARM Arithmetic [95]}
\textbf{Gain:} \texttt{qarm\_gain}. \textbf{Metadata:} \texttt{modulus},
\texttt{orbit\_size}, \texttt{generator}.

\subsection{Symbolic Search [96]}
\textbf{Gain:} \texttt{sym\_gain}. \textbf{Metadata:} \texttt{beam\_width},
\texttt{search\_depth}, \texttt{rule\_count}. \textbf{Example:} Substrate
$(3,5,8,13)$, $\hqa\approx 0.4235$, $\mathrm{lr}=0.01$, $\mathrm{sym\_gain}=1.1$:
$\kappa\approx0.00466$.

\subsection{Orbit Curvature [97]}
\textbf{Gain:} $\kmin$ derived by full orbit enumeration. \textbf{Description:}
Rather than a single-state snapshot, this family certifies the tightest stability
bottleneck across an entire QA orbit. The validator enumerates the complete orbit
under $T$ modulo $m$, computes $\hqa$ at every state, and reports
$\kmin=\min_t\kappa_t$.

\textbf{Orbit length:} $|\mathcal{O}|=\pi(m)/2$ (Proposition~\ref{prop:half-pisano}).
For $m=9$: $\pi(9)=24$, orbit length 12, 72 starting pairs. \textbf{Example:}
Orbit start $(b_0,e_0)=(1,2)$, modulus 9, $\mathrm{lr}=0.5$, $\mathrm{gain}=1.0$:
$\kmin\approx0.1077$ (bottleneck at $(8,1,9,1)$), orbit length 12.

\subsection{GNN Spectral Gain [98] and Attention Spectral Gain [99]}
\textbf{Gain (GNN):} $\sigma_{\max}(W)$ via power iteration on $W^\top W$.
\textbf{Gain (Attention):} $\sigma_{\max}(QK^\top/\!\sqrt{d_k})$ via power
iteration. Gate B computes the gain internally and rejects any certificate where
the claimed value disagrees.

\textbf{GNN example:} $W=\bigl(\begin{smallmatrix}0.8&0.2\\0.1&0.6\end{smallmatrix}\bigr)$,
substrate $(3,5,8,13)$, $\hqa\approx0.4235$, $\sigma_{\max}(W)\approx0.8821$,
$\kappa\approx0.00374$.

\textbf{Attention example:} $Q=\bigl(\begin{smallmatrix}0.6&0.4\\0.3&0.7\end{smallmatrix}\bigr)$,
$K=\bigl(\begin{smallmatrix}0.5&0.3\\0.2&0.8\end{smallmatrix}\bigr)$, $d_k=2$,
$\sigma_{\max}(QK^\top/\!\sqrt{2})\approx0.6603$, $\kappa\approx0.00280$.

\subsection{Gradient Lipschitz Gain [101]}
\textbf{Gain:} $\min(\|\mathbf{g}\|_2,2.0)$ derived from the submitted gradient.
\textbf{Description:} The L2 norm of the gradient is the natural local Lipschitz
constant of the update step. Gate B recomputes the norm and rejects disagreements.
\textbf{Example:} $\mathbf{g}=[0.3,0.4]$, $\|\mathbf{g}\|_2=0.5$, substrate
$(3,5,8,13)$, $\hqa\approx0.4235$, $\mathrm{lr}=0.01$, $\kappa\approx0.00212$.

\section{The Three-Gate Certificate Pattern}

Any artifact claiming compliance must pass a three-gate validation:

\begin{itemize}
  \item \textbf{Gate A -- Substrate Integrity:} The validator recomputes $\hqa$ from
    raw parameters $(a,b,d,e)$ and compares to the claimed value.
    Failure mode: \texttt{H\_QA\_MISMATCH}.

  \item \textbf{Gate B -- Update Rule Integrity:} The validator checks
    $\mathrm{gain}\in(0,2]$, recomputes $\etaeff$, and verifies the claimed
    $p_{\mathrm{after}}$.
    Failure modes: \texttt{GAIN\_OUT\_OF\_RANGE}, \texttt{UPDATE\_RULE\_MISMATCH}.

  \item \textbf{Gate C -- Curvature Integrity:} Using verified inputs, the validator
    recomputes $\kappa$ and compares.
    Failure mode: \texttt{KAPPA\_MISMATCH}.
\end{itemize}

The failure types are disjoint: a certificate fails at exactly one gate, and the
\texttt{invariant\_diff} field identifies the specific numerical discrepancy.

\section{Comparative Certificate Instantiations}
\label{sec:comparison}

Table~\ref{tab:comparison} shows that once a family is in the certified normal form,
the resulting $\kappa$ depends only on $(\mathrm{lr}, \mathrm{gain}, \hqa)$, not on
the architectural label. We use $\mathrm{lr}=0.01$, $\mathrm{gain}=1.1$.

\begin{table}[h]
\centering
\caption{Cross-family comparison at two substrates ($\mathrm{lr}=0.01$, $\mathrm{gain}=1.1$).}
\label{tab:comparison}
\begin{tabular}{lllrrr}
\toprule
Family & ID & Gain name & $H_{QA}$ & gain & $\kappa$ \\
\midrule
QALM Gradient & [89] & \texttt{gain} & 0.2571 & 1.1 & 0.00283 \\
GNN Aggregation & [93] & \texttt{agg\_gain} & 0.2571 & 1.1 & 0.00283 \\
Attention Layer & [94] & \texttt{attn\_gain} & 0.2571 & 1.1 & 0.00283 \\
QARM Arithmetic & [95] & \texttt{qarm\_gain} & 0.2571 & 1.1 & 0.00283 \\
Symbolic Search & [96] & \texttt{sym\_gain} & 0.2571 & 1.1 & 0.00283 \\
\midrule
QALM Gradient & [89] & \texttt{gain} & 0.4235 & 1.1 & 0.00466 \\
GNN Aggregation & [93] & \texttt{agg\_gain} & 0.4235 & 1.1 & 0.00466 \\
Attention Layer & [94] & \texttt{attn\_gain} & 0.4235 & 1.1 & 0.00466 \\
QARM Arithmetic & [95] & \texttt{qarm\_gain} & 0.4235 & 1.1 & 0.00466 \\
Symbolic Search & [96] & \texttt{sym\_gain} & 0.4235 & 1.1 & 0.00466 \\
\bottomrule
\end{tabular}
\end{table}

For a given substrate, $\kappa$ is identical across all five families despite their
different purposes. This demonstrates that $\kappa$ is sensitive to the underlying
scalar inputs while being independent of architectural specifics.

\section{Theorems}
\label{sec:theorems}

\begin{theorem}[Unified Curvature Normal Form]
\label{thm:normal-form}
Let a certified family provide a one-step update
\[
  p_{\mathrm{after}} = p_{\mathrm{before}} - \etaeff\cdot\mathrm{grad},
  \qquad \etaeff=\mathrm{lr}\cdot\mathrm{gain}\cdot\hqa,
\]
where $\hqa$ is computed from a four-parameter substrate $(a,b,d,e)$ by the
definitions in Section~\ref{sec:substrate}, $\mathrm{lr}\geq 0$, $\mathrm{gain}\geq 0$.
Then the associated one-step curvature score $\kappa=1-|1-\etaeff|$ is positive if
and only if $0<\etaeff<2$. In particular, any certified family whose validator
enforces this normal form inherits the same scalar one-step stability criterion,
regardless of domain-specific interpretation of the gain witness.
\end{theorem}

\begin{remark}
This theorem is a statement about a shared certified local update form. It does not
assert global convergence or semantic identity across architecture classes. For QA
orbit families, Theorems~\ref{thm:descent-quad} and~\ref{thm:descent-pl} extend
this to exact multi-step bounds.
\end{remark}

\subsection{Finite-Orbit Descent Theorem (Quadratic Loss)}

The finite structure of the QA orbit enables a multi-step descent result.

\begin{theorem}[Finite-Orbit Descent, Quadratic Loss]
\label{thm:descent-quad}
Let $L(w)=\tfrac{1}{2}w^2$ be a scalar quadratic loss. Let
$\mathcal{O}=\{s_0,\ldots,s_{L-1}\}$ be a QA cosmos orbit of length
$L=\pi(m)/2$ with per-step effective rates
$\etaeff^{(t)}=\mathrm{lr}\cdot\mathrm{gain}\cdot\hqa(s_t)$ and curvature scores
$\kappa_t=1-|1-\etaeff^{(t)}|$. Define the orbit contraction factor
\[
  \rhoO := \prod_{t=0}^{L-1}(1-\kappa_t)^2.
\]
If $\kmin(\mathcal{O}):=\min_t\kappa_t>0$, then:
\begin{enumerate}[(i)]
  \item $L_{t+L}(w)\leq\rhoO\cdot L_t(w)$ for every initial $w$;
  \item $\rhoO\leq(1-\kmin)^{2L}<1$;
  \item $\rhoO$ is a computable exact quantity from the orbit alone, with no
    stochastic approximation.
\end{enumerate}
\end{theorem}

\begin{proof}
For the scalar quadratic, $w_{t+1}=(1-\etaeff^{(t)})w_t$, so
$L_{t+1}=(1-\etaeff^{(t)})^2 L_t$. Since
$|1-\etaeff^{(t)}|=1-\kappa_t$, we have $L_{t+1}=(1-\kappa_t)^2 L_t$.
Composing over $L$ steps: $L_{t+L}=\rhoO\cdot L_t$.
For (ii): $(1-\kappa_t)^2\leq(1-\kmin)^2$ for all $t$, so
$\rho\leq(1-\kmin)^{2L}<1$ since $\kmin>0$.
For (iii): each $\kappa_t$ is computable from $s_t$ via Section~\ref{sec:substrate},
and $L=\pi(m)/2$ is finite by Proposition~\ref{prop:half-pisano}.
\end{proof}

\begin{corollary}[Multi-orbit convergence]
After $k$ full orbits, $L_{kL}\leq\rhoO^k\cdot L_0\to 0$ geometrically with rate
$\log(1/\rhoO)$ per orbit.
\end{corollary}

\paragraph{Key identity.}
The per-step loss contraction factor is $(1-\kappa_t)^2$: higher $\kappa$ directly
and exactly governs loss decrease. This gives the first-principles explanation of
the empirical correlation $r(\text{mean}\,\kappa,\,\text{final loss})=-0.843$
reported in Section~\ref{sec:empirical}.

\paragraph{Numerical example ($m=9$, $\mathrm{lr}=0.5$, $\mathrm{gain}=1$).}
The mod-9 cosmos orbit has $L=12$ states. Numerical evaluation yields
$\kmin=0.1077$ (bottleneck at state $(8,1,9,1)$) and
\[
  \rhoO = \prod_{t=0}^{11}(1-\kappa_t)^2 = 0.001582,
\]
so every full orbit reduces loss by a factor of 632. The orbit-invariant $\rhoO$ is
the same for all 72 starting pairs in the cosmos group (they all traverse the same
12-cycle). After 10 orbits: $L_{10L}\leq(0.001582)^{10}\cdot L_0\approx 10^{-28}\cdot L_0$.

\subsection{Extension to PL-Conditioned Losses}
\label{sec:pl}

The scalar quadratic is a special case ($\beta=\mu$) of a broader class.

\begin{theorem}[Finite-Orbit Descent, PL Loss]
\label{thm:descent-pl}
Let $L:\mathbb{R}^d\to\mathbb{R}$ be $\beta$-smooth and satisfy the
Polyak-\L{}ojasiewicz (PL) condition with constant $\mu>0$:
$\|\nabla L(w)\|^2\geq 2\mu(L(w)-L^*)$ for all $w$,
where $L^*=\inf_w L(w)$. Let $\mathcal{O}$ be a QA cosmos orbit with
$\kmin(\mathcal{O})>0$ and $\etaeff^{(t)}<2/\beta$ for all $t$. Define
\[
  \rhoPL := \prod_{t=0}^{L-1}\!\left(1-2\mu\,\etaeff^{(t)}\!\left(1-\tfrac{\beta}{2}\etaeff^{(t)}\right)\right).
\]
Then $0\leq\rhoPL<1$ and
$L(w_{t+L})-L^*\leq\rhoPL\cdot(L(w_t)-L^*)$.
\end{theorem}

\begin{proof}
At each step, the $\beta$-smooth descent lemma gives
$L(w_{t+1})\leq L(w_t)-\etaeff^{(t)}(1-\tfrac{\beta}{2}\etaeff^{(t)})\|\nabla L(w_t)\|^2$.
Applying the PL condition:
$L(w_{t+1})-L^*\leq\bigl(1-2\mu\etaeff^{(t)}(1-\tfrac{\beta}{2}\etaeff^{(t)})\bigr)(L(w_t)-L^*)$.
Since $\etaeff^{(t)}\in(0,2/\beta)$, each factor is in $[0,1)$.
Composing over $L$ steps gives $\rhoPL<1$.
\end{proof}

\paragraph{Reduction to Theorem~\ref{thm:descent-quad}.}
For quadratic loss with Hessian $\lambda I$ (so $\beta=\mu=\lambda$), the per-step
factor becomes $1-2\lambda\eta(1-\tfrac{\lambda}{2}\eta)=(1-\lambda\eta)^2=(1-\kappa_t)^2$,
so $\rhoPL=\rhoO$. This identity holds exactly for all $\etaeff\in(0,2)$ (verified
numerically). Theorem~\ref{thm:descent-quad} is the canonical case of
Theorem~\ref{thm:descent-pl}.

\paragraph{Condition number dependence.}
For the mod-9 cosmos orbit at $\mathrm{lr}=0.5$: $\rhoPL=0.0016$ (condition 1),
$0.61$ (condition 10), $0.95$ (condition 100). The orbit still guarantees descent for
any finite condition number, with rate degrading gracefully.

\section{Limitations}

For general nonconvex losses, this framework certifies only a local one-step normal
form and does not prove multi-step convergence, generalization, or robustness under
distribution shift. Theorems~\ref{thm:descent-quad} and~\ref{thm:descent-pl} give
exact multi-step guarantees for scalar quadratic and $\beta$-smooth $\mu$-PL losses
respectively; extension to nonconvex losses (beyond PL) requires additional
assumptions and remains open.

In families [89]--[96], the gain is supplied as a free scalar witness; the metadata
fields (\texttt{n\_nodes}, \texttt{d\_model}, \texttt{orbit\_size}, etc.) are
recorded for auditing but do not constrain that scalar. Accordingly, those families
should be read as reusable machine-checkable local certification patterns. Families
[98], [99], and [101] close this gap: gain is derived from $\sigma_{\max}(W)$,
$\sigma_{\max}(QK^\top/\!\sqrt{d_k})$, and $\|\mathbf{g}\|_2$ respectively, and
cannot be freely adjusted. Extending derived-gain derivations to the remaining
families (QARM, symbolic search) is left to future work.

\section{Conclusion}

This paper introduced a certificate-normal-form framework grounded in the arithmetic
of $\Qfi$: the QA map $T=Q^2$ acts as multiplication by $\varphi^2$ in $\Zfi$, and
the harmonic index $\hqa$ decomposes as the normalised trace of a natural
cross-coupling operator on the QA step matrix. Building on this foundation, we
derived a universal curvature score $\kappa$ and certified its one-step stability
form across nine architecture families.

For QA orbit families, Theorems~\ref{thm:descent-quad} and~\ref{thm:descent-pl}
provide exact multi-step descent guarantees: if $\kmin(\mathcal{O})>0$, one full
orbit of length $\pi(m)/2$ reduces loss by exactly computable factor $\rhoO<1$ (for
quadratic loss) or $\rhoPL<1$ (for any $\beta$-smooth $\mu$-PL loss). For the
mod-9 cosmos orbit at $\mathrm{lr}=0.5$, $\rhoO=0.001582$ (factor-632 per orbit).
These results are enabled uniquely by the finite enumerability of QA orbits---a
structural property with no direct analogue in standard SGD analysis.

Future work: (1) extend to nonconvex losses via Lyapunov arguments; (2) derive gain
from native objects in QARM and symbolic search families; (3) explore $\kappa$ as an
active regularisation term in training; (4) develop a standalone mathematical note
on the $\Zfi$/Half-Pisano/norm-invariant structure.

\section{Empirical Validation}
\label{sec:empirical}

We validate the $\kappa$ framework with two controlled experiments on synthetic
binary classification (800 samples, 20 features, two Gaussian classes). The model is
a two-layer MLP ($20\to32$ ReLU$\to1$ sigmoid) trained with binary cross-entropy,
pure NumPy, no GPU.

\paragraph{Setup.}
\[
  p_{\mathrm{after}} = p_{\mathrm{before}}
    - \mathrm{lr}\cdot\mathrm{gain}\cdot\hqa\cdot\nabla\mathcal{L},
  \qquad
  \kappa = 1-|1-\mathrm{lr}\cdot\mathrm{gain}\cdot\hqa|.
\]

\paragraph{Experiment 1 -- $H_{QA}$ sweep.}
Seven substrates spanning $\hqa\in[0.20,0.71]$ (mod-9, gain $=\min(\|\mathbf{g}\|_2,2)$,
$\mathrm{lr}=0.1$, 300 epochs) plus plain-SGD baseline.

\begin{table}[h]
\centering
\caption{Experiment 1: H\_QA sweep results.}
\begin{tabular}{lrrr}
\toprule
Substrate & $H_{QA}$ & mean $\kappa$ & Final loss \\
\midrule
(2,8) & 0.201 & 0.0035 & 0.0434 \\
(4,7) & 0.305 & 0.0043 & 0.0395 \\
(1,4) & 0.356 & 0.0052 & 0.0407 \\
(5,3) & 0.471 & 0.0056 & 0.0289 \\
(9,8) & 0.529 & 0.0065 & 0.0279 \\
(3,5) & 0.594 & 0.0068 & 0.0254 \\
(1,5) & 0.715 & 0.0077 & 0.0231 \\
\bottomrule
\end{tabular}
\end{table}

The ordering is monotone: higher $\hqa\to$ higher $\kappa\to$ lower final loss.
Pearson correlations: $r(\text{mean}\,\kappa,\,\text{final loss})=-0.843$,
$r(\hqa,\,\text{final accuracy})=+0.769$. Plain SGD achieves lower absolute loss
because gain $=\|\mathbf{g}\|_2\ll 1$ in the QA conditions shrinks the effective
step by 20--150$\times$; within the QA family the $\kappa$ ordering is preserved
and strongly predictive.

\paragraph{Experiment 2A -- $\eta_{\mathrm{eff}}$ sweep.}
Fixed substrate $(9,8)$ ($\hqa=0.529$), gain $=1$, lr $=\etaeff/\hqa$ exact.

\begin{table}[h]
\centering
\caption{Experiment 2A: $\eta_{\mathrm{eff}}$ sweep (substrate (9,8)).}
\begin{tabular}{rrr}
\toprule
$\etaeff$ & $\kappa$ & Final loss \\
\midrule
0.05 & 0.05 & $1.77\times10^{-3}$ \\
0.25 & 0.25 & $2.32\times10^{-4}$ \\
0.50 & 0.50 & $1.05\times10^{-4}$ \\
1.00 & 1.00 & $4.6\times10^{-5}$ \\
1.50 & 0.50 & $1.1\times10^{-5}$ \\
1.90 & 0.10 & $1.9\times10^{-5}$ \\
\bottomrule
\end{tabular}
\end{table}

Loss decreases monotonically from $\etaeff=0.05$ to $1.50$, then stabilises.
Pearson $r(\kappa,\,\text{final loss})=-0.53$.

\paragraph{Experiment 2B -- substrate equalization.}
Setting lr $=1/\hqa$ for each of seven substrates forces $\kappa=1$ uniformly.
All seven converge in epoch 1; loss standard deviation $1.9\times10^{-5}$
(range $[7\times10^{-6},7\times10^{-5}]$). This confirms that $\hqa$ governs
convergence only through $\etaeff$; once $\etaeff$ is equalized, substrate
identity has negligible effect.

\paragraph{Summary.}
Three claims supported jointly: (i) within the QA update family, mean $\kappa$ is
a strong predictor of convergence quality ($r=-0.843$); (ii) the critical design
parameter is $\etaeff$, with $\kappa$ characterizing proximity to the stability
boundary; (iii) once $\etaeff$ is equalized, convergence is equalized.

\section{Appendix: Mathematical Proofs}
\label{sec:appendix}
\setcounter{theorem}{0}

\begin{proposition}[Half-Pisano Orbit Length]
\label{prop:half-pisano}
For every $m\geq 3$, the QA update map
$T:(b,e)\mapsto(b{+}e\bmod_m^*,\,b{+}2e\bmod_m^*)$ satisfies $T=Q^2$ and has
maximal orbit length $\pi(m)/2$, where $\pi(m)$ is even for all $m\geq 3$.
A state $(b_0,e_0)\in\{1,\ldots,m\}^2$ is \emph{primitive} if its orbit under $T$
attains this maximal length.
\end{proposition}

\begin{proof}[Proof sketch]
\textit{Step 1 -- Matrix identity.}
\[
  T\begin{pmatrix}b\\e\end{pmatrix}
  = \begin{pmatrix}1&1\\1&2\end{pmatrix}\begin{pmatrix}b\\e\end{pmatrix}
  = Q^2\begin{pmatrix}b\\e\end{pmatrix}\pmod{m}.
\]

\textit{Step 2 -- Pisano period is even for $m\geq 3$.}
Wall~\cite{wall1960} shows that for $m\geq 3$, $\pi(m)\equiv 0\pmod{2}$.

\textit{Step 3 -- Order of $Q^2$.}
$\operatorname{ord}(Q^2)=\pi(m)/\gcd(2,\pi(m))=\pi(m)/2$.

\textit{Step 4 -- Orbit length.}
The orbit length of $(b_0,e_0)$ under $T=Q^2$ divides $\pi(m)/2$. A state is
primitive if its orbit length equals $\pi(m)/2$; non-primitive states yield shorter
orbits.
\end{proof}

\begin{proposition}[Quadratic Norm Invariant and $\mathbb{Q}(\sqrt{5})$ Structure]
\label{prop:norm-invariant}
The QA map $T$ is the linear representation of multiplication by $\varphi^2$ in
$\Zfi$. The quadratic form $f(b,e)=b^2+be-e^2$ is the algebraic norm
$N(b+e\varphi)$ in $\Qfi$, and is invariant under $T$. The mod-$m$ value of $f$
classifies QA orbits.
\end{proposition}

\begin{proof}
\textit{Step 1 -- Algebraic identification.}
Multiplication by $\varphi$ has matrix $Q$ (since
$\varphi(b+e\varphi)=b\varphi+e\varphi^2=e+(b+e)\varphi$). Therefore $T=Q^2$
represents multiplication by $\varphi^2$.

\textit{Step 2 -- Norm identification.}
$N(b+e\varphi)=(b+e\varphi)(b+e\bar\varphi)=b^2+be(\varphi+\bar\varphi)+e^2\varphi\bar\varphi
=b^2+be-e^2=f(b,e)$, using $\varphi+\bar\varphi=1$ and $\varphi\bar\varphi=-1$.

\textit{Step 3 -- Invariance.}
Since $N(\varphi^2)=N(\varphi)^2=(-1)^2=1$, multiplication by $\varphi^2$ preserves
the norm: $f(T(b,e))=f(b,e)$.
\end{proof}

\paragraph{Geometric identity.}
$\det\bigl(\begin{smallmatrix}b&d\\e&a\end{smallmatrix}\bigr)
=ba-ed=b(b+2e)-e(b+e)=b^2+be-e^2=f(b,e)$. Thus $f(b,e)$ is the determinant of
the matrix whose columns are the state before and after one QA step.

\begin{corollary}[Orbit classification, $m=9$]
The orbit type of $(b_0,e_0)$ is determined by the 3-adic valuation of
$N(b_0+e_0\varphi)$:

\begin{center}
\begin{tabular}{lll}
\toprule
$v_3(f(b_0,e_0))$ & Orbit length & Orbit type \\
\midrule
$\geq 4$ & 1 & Singularity $(9,9)$ \\
$2$ or $3$ & 4 & Satellite \\
$0$ & 12 & Cosmos (primitive) \\
\bottomrule
\end{tabular}
\end{center}
Verified for all 81 states of $\{1,\ldots,9\}^2$.
\end{corollary}

\paragraph{Fibonacci projective flow.}
The projective ratio $z_t=e_t/(b_t+e_t)$ evolves as $z_{t+1}=(1+z_t)/(2+z_t)$
under $T$, with unique fixed point $z^*=(\sqrt{5}-1)/2=1/\varphi\approx0.618$.

\paragraph{$H_{QA}$ decomposition.}
Setting $d=b+e$ and $z=e/d$:
\[
  \hraw = \frac{1-z^2}{4(1+z^2)} + \frac{e}{8}.
\]
The first term equals $\cos(2\theta)/4$ where $\theta=\arctan(z)$; the second,
$e/8=\sigma_{od}/4$, follows from the Fibonacci identity $a+b=2d$. Writing
$M=\bigl(\begin{smallmatrix}b&e\\d&a\end{smallmatrix}\bigr)$:
\[
  4\hraw = \sigma_d(M)+\sigma_{od}(M),\quad
  \sigma_d=\frac{M_{11}M_{22}}{M_{12}^2+M_{21}^2}\;(\text{degree-0}),\quad
  \sigma_{od}=\frac{M_{12}M_{21}}{M_{11}+M_{22}}\;(\text{degree-1}).
\]
The formula satisfies transpose symmetry $\sigma(M)=\sigma(M^\top)$ and bilinear
identity $4\hraw(e^2+d^2)(b+a)=ba(b+a)+ed(e^2+d^2)$ (verified for all seven
experimental substrates). Defining $\Lambda=\mathrm{diag}(\sigma_d,\sigma_{od})$:
$4\hraw=\mathrm{tr}(\Lambda)$ and $2\hraw=R(\Lambda,(1,1)^\top/\sqrt{2})$ (the
Rayleigh quotient at the equal-weight vector). The factor $1/4=\mathrm{tr}(\Lambda)/(2n)|_{n=2}$
is the normalised-trace scaling of the $2\times 2$ operator.

\begin{theorem}[Orbit Contraction]
Let $\mathcal{O}=\{s_0,\ldots,s_{L-1}\}$ be a finite QA orbit of length $L=\pi(m)/2$
with per-step effective rate $\etaeff^{(t)}=\mathrm{lr}\cdot\mathrm{gain}\cdot\hqa^{(t)}$.
Define $\kmin(\mathcal{O})=\min_t(1-|1-\etaeff^{(t)}|)$.
If $\kmin(\mathcal{O})>0$, then $0<\etaeff^{(t)}<2$ for every $t\in\{0,\ldots,L-1\}$.
\end{theorem}
\begin{proof}
$\kmin>0\Leftrightarrow|1-\etaeff^{(t)}|<1\Leftrightarrow 0<\etaeff^{(t)}<2$ for all $t$.
Since the orbit is finite and enumerable, this condition can be checked exactly.
\end{proof}

\bibliographystyle{plain}
\begin{thebibliography}{9}

\bibitem{vaswani2017attention}
A. Vaswani et al.
\newblock Attention is all you need.
\newblock In \textit{Advances in Neural Information Processing Systems}, 2017.

\bibitem{wu2020comprehensive}
Z. Wu, S. Pan, F. Chen, G. Long, C. Zhang, and P. S. Yu.
\newblock A comprehensive survey on graph neural networks.
\newblock \textit{IEEE Trans. Neural Netw. Learn. Syst.}, 32(1):4--24, 2020.

\bibitem{kingma2014adam}
D. P. Kingma and J. Ba.
\newblock Adam: A method for stochastic optimization.
\newblock \textit{arXiv:1412.6980}, 2014.

\bibitem{bottou2012sgd}
L. Bottou.
\newblock Stochastic gradient descent tricks.
\newblock In \textit{Neural Networks: Tricks of the Trade}, pp.~421--436. Springer, 2012.

\bibitem{kipf2016gcn}
T. N. Kipf and M. Welling.
\newblock Semi-supervised classification with graph convolutional networks.
\newblock \textit{arXiv:1609.02907}, 2016.

\bibitem{ruder2016overview}
S. Ruder.
\newblock An overview of gradient descent optimization algorithms.
\newblock \textit{arXiv:1609.04747}, 2016.

\bibitem{silver2016go}
D. Silver et al.
\newblock Mastering the game of Go with deep neural networks and tree search.
\newblock \textit{Nature}, 529(7587):484--489, 2016.

\bibitem{wolf2019transformers}
T. Wolf et al.
\newblock HuggingFace's Transformers: State-of-the-art natural language processing.
\newblock \textit{arXiv:1910.03771}, 2019.

\bibitem{wall1960}
D. D. Wall.
\newblock Fibonacci primitive roots and the period of the Fibonacci sequence modulo a prime.
\newblock \textit{Am. Math. Monthly}, 67(6):525--532, 1960.

\end{thebibliography}

\end{document}
